\documentclass[tikz,border=6pt]{standalone}
\usepackage{ctex}
\usepackage{amsmath}
\usetikzlibrary{calc, arrows.meta}

\begin{document}
\begin{tikzpicture}[scale=1.1, thick]

% 焦点弦示意图：过焦点的弦 AB 三大性质
% 抛物线 y^2 = 4x (p=2)，焦点 F(1,0)，准线 x=-1

% 坐标轴
\draw[->, thin, gray] (-2.0, 0) -- (6.0, 0) node[right, font=\small] {$x$};
\draw[->, thin, gray] (0, -3.8) -- (0, 3.8) node[above, font=\small] {$y$};
\node[font=\footnotesize, below right] at (0,0) {$O$};

% 抛物线
\draw[black, thick, domain=-3.3:3.3, smooth, samples=100]
  plot ({(\x)^2/4}, \x);

% 准线 x = -1
\draw[red, thin] (-1, -3.6) -- (-1, 3.6);
\node[red, font=\footnotesize, above] at (-1, 3.6) {准线};

% 焦点 F
\fill[black] (1, 0) circle [radius=2.5pt];
\node[below left, font=\small] at (1, 0) {$F$};

% 焦点弦：斜率 k=1，直线 y = x-1，交 y^2=4x
% y^2=4(y+1) => y^2-4y-4=0 => y=(4±sqrt(32))/2=2±2√2
% A: y=2+2√2≈4.828 => x=(y+1)=y/k+1... 用参数化：A(t_A=2+2√2), B(t_B=2-2√2)
% t_A=2+2*1.414=4.828, x_A=t_A^2/4=5.827
% t_B=2-2*1.414=-0.828, x_B=t_B^2/4=0.171
\coordinate (A) at ({(2+2*1.4142)^2/4}, {2+2*1.4142}); % ≈(5.83, 4.83)
\coordinate (B) at ({(2-2*1.4142)^2/4}, {2-2*1.4142}); % ≈(0.17, -0.83)

% 弦 AB（蓝色）
\draw[blue, thick] (A) -- (B);

% 标注 A, B
\fill[blue] (A) circle [radius=2.5pt];
\fill[blue] (B) circle [radius=2.5pt];
\node[above right, font=\small, blue] at (A) {$A(x_1,y_1)$};
\node[below right, font=\small, blue] at (B) {$B(x_2,y_2)$};

% AF, BF 标注
\draw[blue, dashed] (A) -- (1,0)
  node[midway, left, font=\footnotesize] {$|AF|$};
\draw[blue, dashed] (B) -- (1,0)
  node[midway, below right, font=\footnotesize] {$|BF|$};

% A, B 在准线上的射影
\coordinate (PA) at (-1, {2+2*1.4142});
\coordinate (PB) at (-1, {2-2*1.4142});
\draw[gray, dashed, thin] (A) -- (PA);
\draw[gray, dashed, thin] (B) -- (PB);
\fill[gray] (PA) circle [radius=1.5pt];
\fill[gray] (PB) circle [radius=1.5pt];
\node[left, font=\footnotesize, gray] at (PA) {$A'$};
\node[left, font=\footnotesize, gray] at (PB) {$B'$};

% 三大性质文本框
\node[font=\footnotesize, align=left,
      draw=blue!50, rounded corners=3pt, fill=blue!5,
      inner sep=5pt] at (4.5, -2.5)
  {焦点弦三大性质：\\
   $|AB|=x_1+x_2+p$\\[2pt]
   $\dfrac{1}{|AF|}+\dfrac{1}{|BF|}=\dfrac{2}{p}$\\[2pt]
   $|AF|\cdot|BF|=\dfrac{p^2}{...}$（见正文）};

\end{tikzpicture}
\end{document}
